%======================================================================
\subsection{BW-PDM-08 --- Plan: migrating from Fusion (24 Oct 2026, approved 24 Oct)}\label{sec:pdm08-plan}
%======================================================================

\eref{PDM-08} moves the Fusion designs into the vault, numbered and at the right revision, and turns the FC-T01 switch
criteria into the checklist that closes M6. It comes last in the track (it depends on the pilot, \eref{PDM-07}), but the
facts below decide what the export from Fusion must contain, and that is worth knowing before the pilot exports
anything. The acceptance stays as written: \emph{all LE parts imported, numbered and released; FC-T01 criteria met.}

%----------------------------------------------------------------------
\subsubsection*{1. What Fusion can give, and what survives}

\begin{center}
\small
\begin{tabularx}{\linewidth}{@{}>{\raggedright\arraybackslash}p{2.3cm}>{\hsize=1.35\hsize}X>{\raggedright\arraybackslash}p{2.7cm}>{\hsize=0.65\hsize}X@{}}
\toprule
Export & Holds & FreeCAD reads & Use in the migration\\
\midrule
\file{.f3d} & One design: history, sketches, parameters, joints, appearances, materials & no (Autodesk's format) &
  archive only, attached to the item\\
\file{.f3z} & A design with the designs it references & no & archive only, attached to the assembly\\
STEP (AP214, \verify{} AP242) & B-rep solids; one \code{PRODUCT} per component and one occurrence (NAUO)
  per instance, with placements, names, colours. \textbf{No history,
  parameters, sketches or joints.} & yes (\code{Import}) & the geometry\\
STL, 3MF, OBJ & meshes & yes, as meshes & print files only; not migrated as models\\
Drawing PDF & the sheet as vector graphics & no & attached as the released drawing of a migrated revision\\
Drawing DWG, DXF & 2D geometry and text, not associative & yes (DXF; DWG with a converter) & not needed if the PDF is kept\\
Fusion API (Python) & per component: name, part number, description, material, physical properties (mass, volume,
  area, centre of mass), occurrences with transforms; STEP per component & --- & the manifest that makes the import
  exact (section~3)\\
\bottomrule
\end{tabularx}
\end{center}

\begin{center}
\small
\begin{tabularx}{\linewidth}{@{}XXX@{}}
\toprule
Survives the move & Rebuilt in FreeCAD & Lost (kept only in the archive)\\
\midrule
Solid geometry, exact B-rep & joints (\eref{ASM-03}) & feature history, sketches\\
Assembly structure, quantities, placements & materials, from the manifest & user parameters\\
Names, part numbers (manifest) & drawings, when a part next changes & drawing associativity\\
Colours (AP214) & sheet-metal flat patterns (\eref{MOD-06}) & CAM set-ups\\
\bottomrule
\end{tabularx}
\end{center}

%----------------------------------------------------------------------
\subsubsection*{2. What the spike showed (24 Oct)}

A STEP file of a two-level assembly was written by this build (\file{pdm-plans/spike\_step.py}) with Fusion-style names
(\code{LE-02001-A.1 Pedestal Assy v7} holding \code{LE-01001-A.1 Base plate v12} and two instances of
\code{LE-01010-A.1 Lower vertical supports v3}), then read back three ways.
\begin{center}
\small
\begin{tabularx}{\linewidth}{@{}>{\hsize=1.15\hsize}X>{\hsize=0.85\hsize}X@{}}
\toprule
Check & Result\\
\midrule
Product tree read from the STEP text (50 lines of Python, no CAD kernel, \file{pdm-plans/step\_tree.py}) & 3
  products, support $\times2$: the structure and quantities are in the file\\
\code{Import.insert} in \code{FreeCADCmd}, default options & works; the two support instances come back as
  \emph{two separate copies}, not links\\
\quad labels of the copies & \code{\ldots supports v3} and \code{\ldots supports v001}: FreeCAD's unique-label rule
  rewrote the trailing ``3'' of Fusion's version suffix\\
\quad with \code{useLinkGroup=True} (instances as links) & \textbf{crash} (segmentation fault), twice\\
\quad with one-document-per-part modes (\code{mode=1}, \code{3}, \code{merge=False}) & ignored headless: one document,
  copies (\verify{} in the GUI)\\
Vault import (\file{vault\_setup.py}) & reads \file{.FCStd} only; \code{LE-01001-A.1} is not a vault number, so a file
  carrying it gets a new one\\
Vault revision letters & \code{S}, \code{X}, \code{Z} refused; LE-STD-0001 allows them\\
\bottomrule
\end{tabularx}
\end{center}

\paragraph{What follows.} Matching FreeCAD objects to Fusion components by label is unsafe (labels are rewritten), and
FreeCAD's own assembly import either copies instances or, headless, crashes when asked to keep them as links. So the
migration does not rely on importing a whole assembly: \textbf{each component is exported as its own STEP file and the
structure comes from a manifest}, from which the vault builds the parts and then the assemblies with links, exactly
as \eref{ASM-01} wants them. The STEP text reader stays as a cross-check of the manifest's quantities. The crash is
to be reported upstream with the spike file.

%----------------------------------------------------------------------
\subsubsection*{3. The pipeline}

\begin{figure}[!htbp]
\centering
\begin{tikzpicture}[>=Stealth,node distance=5mm and 5mm,
  b/.style={draw=bwblue!70,fill=bwblue!5,rounded corners=2pt,font=\scriptsize,align=center,minimum height=1.05cm,text width=2.2cm},
  f/.style={b,draw=bworange!80,fill=bworange!8},
  u/.style={b,draw=bwgreen!80,fill=bwgreen!8},
  t/.style={font=\tiny\color{bwgrey},align=center,text width=2.6cm}]
\node[f] (fu) {Fusion project\\\emph{Leonardo / v0\_proto}};
\node[f,right=of fu] (ex) {Export script\\(Fusion API)};
\node[b,right=of ex] (st) {Staging folder\\STEP per component, manifest, PDFs, \texttt{.f3d}};
\node[b,right=of st] (ck) {Check: register, numbers, revisions, duplicates};
\node[u,below=9mm of ck] (cs) {\texttt{migration.csv}\\you review it};
\node[b,left=of cs] (cv) {Convert: one \texttt{.FCStd} per part, assemblies with links};
\node[b,left=of cv] (vf) {Verify: mass properties, flattened BOM};
\node[b,left=of vf] (im) {Vault import: numbers, revisions, provenance, attachments};
\node[u,below=9mm of im] (fc) {FC-T01 checklist};
\draw[->] (fu) -- (ex); \draw[->] (ex) -- (st); \draw[->] (st) -- (ck); \draw[->] (ck) -- (cs);
\draw[->] (cs) -- (cv); \draw[->] (cv) -- (vf); \draw[->] (vf) -- (im); \draw[->] (im) -- (fc);
\node[t,above=1mm of ex] {runs in Fusion, by you};
\node[t,above=1mm of ck] {dry run: writes nothing};
\node[t,below=1mm of vf,text width=2.6cm] {fails the item if\\$|\Delta V|/V > 10^{-4}$};
\node[t,right=4mm of fc,text width=6.2cm,align=left] {Lost on the way, kept in the attached \texttt{.f3d}: history,
  sketches, parameters, joints. Rebuilt: joints, materials, drawings when a part next changes.};
\end{tikzpicture}
\caption{Migration from Fusion. Orange runs in Fusion, blue in Heuretes FreeCAD, green are the two points where you decide.
Nothing is written to the vault before the reviewed \file{migration.csv} says so.}
\label{fig:pdm08-pipeline}
\end{figure}

\begin{itemize}[nosep]
  \item \textbf{Export script} (Python, Fusion API, run by you in Fusion). For every component in the chosen designs:
        its STEP file, and a line in \file{manifest.json} with the component's name, part number and description fields,
        material, physical properties, the design's Fusion version, and every occurrence with its parent and transform.
        Each drawing as PDF, each design as \file{.f3d} (\file{.f3z} for assemblies). The Fusion API works in
        centimetres, so lengths are converted once, in the script.
  \item \textbf{Check} (Heuretes FreeCAD, FreeCAD-free, like \code{scan\_for\_import}). Reads the manifest and the register,
        proposes a vault number and revision for each component (section~4), finds repeats (the same component in two
        designs), and writes \file{migration.csv}. Nothing is reserved yet.
  \item \textbf{Review.} You open \file{migration.csv} in a spreadsheet, fix a number or a revision, untick a row; the
        tool reads it back, as Vault Settings does with its CSV.
  \item \textbf{Convert} (\code{FreeCADCmd}). Each part's STEP becomes \file{<P/N>.FCStd} with an \code{App::Part},
        its solid and the part properties; each assembly becomes an Assembly with one \code{App::Link} per occurrence
        (placement from the manifest), linking the part files by P/N, children first, as \code{run\_import} already
        orders them.
  \item \textbf{Verify} (section~4), then \textbf{import}: the existing import (\code{run\_import}) learns to take the
        converted files with numbers fixed by the CSV, records provenance in the check-in, and attaches the PDF and
        \file{.f3d} as write-once objects.
\end{itemize}

%----------------------------------------------------------------------
\subsubsection*{4. Theory: mapping numbers and revisions, and checking geometry}

\paragraph{Numbers.} A Fusion component $f$ is mapped to a vault number by the first rule that applies:
\[
  \varphi(f) =
  \begin{cases}
    \text{the LE ID in its part-number field or name}, & \text{matching } \texttt{LE-}(\mathit{BB})(\mathit{SSS}) \\
    \text{the register row with the same title}, & \text{if exactly one}\\
    \text{the next free number in its band}, & \text{never a blocked one}\\
  \end{cases}
\]
with the band from the register class (part, assembly, tooling, harness, electronic, reference, COTS). Two components
that map to the same number are a conflict for you to resolve in the CSV, never merged silently. Scratch (\code{ZZ})
components are skipped. The number form itself follows question~2 of \S\ref{sec:pdm07-plan}.

\paragraph{Revisions.} Fusion's own version numbers (v7, v12) are saves, not revisions; they go into the provenance
only. The LE revision $\ell$ on the register row maps to a vault state and revision:
\[
  \psi(\ell) =
  \begin{cases}
    (\text{WIP},\ 0n) & \ell = \text{A.}n \quad\text{(never released)}\\
    (\text{Released},\ X) & \ell = X\\
    (\text{Released } X^{-},\ \text{then WIP } X \text{ under a migration ECO}) & \ell = X.n,\ X \ne \text{A}
  \end{cases}
\]
where $X^{-}$ is the letter before $X$. The vault's lifecycle cannot create a released revision without a review, so the
import writes a new journal event, \code{import\_release}, which records the revision as released \emph{by migration}
with its provenance (Fusion design, version, export date, SHA-256 of the STEP and the \file{.f3d}) instead of a
checklist. Earlier letters are recorded as superseded, without a file unless an archive of them exists (question~4).
With the display rule of \S\ref{sec:pdm07-plan} (question~3), \code{0}$n$ prints as \code{A.}$n$, so nobody sees the
mapping.

\paragraph{Geometry: a checksum that survives conversion.} A STEP file's bytes change with every export (it carries a
timestamp), so its SHA-256 cannot show that a part arrived intact. Mass properties can. For each part the manifest gives
Fusion's volume $V$, area $A$ and centre of mass $\mathbf{c}$; FreeCAD computes the same from the converted solid. The part
passes when
\[
  \frac{|V' - V|}{V} < 10^{-4},\qquad \frac{|A' - A|}{A} < 10^{-4},\qquad \lVert\mathbf{c}' - \mathbf{c}\rVert < 0.01\,\text{mm}.
\]
The B-rep is exact in both kernels, so honest differences come only from their integration tolerances, far below these
bounds; anything above them is a lost face or a body that did not convert. For each assembly the flattened BOM must match: for every leaf
$x$, $Q'_t(x) = Q_t(x)$, with $Q'$ from the vault (\S\ref{sec:pdm06-plan}) and $Q$ from the manifest's occurrences, and
the STEP text reader as a third count.

%----------------------------------------------------------------------
\subsubsection*{5. FC-T01: a proposed switch checklist}

The Ops data marks ``Bruce sets FC-T01 switch criteria'' as done, but the criteria are not recorded (open question
\ref{q:fct01}). Proposed, so there is something to correct:
\begin{center}
\small
\begin{tabularx}{\linewidth}{@{}lX@{}}
\toprule
\# & Criterion\\
\midrule
1 & The parity matrix (Appendix~\ref{app:parity}) is met for every row the LE parts use.\\
2 & \code{LE-02001} rev A was released from FreeCAD (\eref{PDM-07}), and one ECR $\rightarrow$ ECO $\rightarrow$ rev B
  round trip was done on a real part (\eref{PDM-05}).\\
3 & Every issued LE register row is a vault item with the same number and revision; none is unmapped; no blocked number
  was used.\\
4 & Every migrated part passed the geometry check; every assembly's flattened BOM matches Fusion's.\\
5 & Every released drawing exists in the vault (FreeCAD, or the Fusion PDF attached); every part that will change has a
  FreeCAD drawing.\\
6 & Print, CNC and laser files for at least one LE part each were made from FreeCAD (M3).\\
7 & The vault is backed up off this Mac, and a restore was tested (\emph{Move vault} verifies by SHA-256).\\
8 & Two weeks of LE work in FreeCAD only, with no fallback to Fusion.\\
9 & The Fusion project is archived into the vault (\file{.f3z} per assembly) and set read-only.\\
\bottomrule
\end{tabularx}
\end{center}

%----------------------------------------------------------------------
\subsubsection*{6. Build order and effort}

\begin{center}
\small
\begin{tabularx}{\linewidth}{@{}lXl@{}}
\toprule
\# & Step & Effort\\
\midrule
1 & Fusion export script: STEP per component, manifest, PDFs, archives (tested on one design with you) & 1.5\,d\\
2 & STEP text reader, manifest check, mapping $\varphi$ and $\psi$, \file{migration.csv} out and back in & 1.5\,d\\
3 & Converter: part files, assemblies with links and placements, children first (\code{FreeCADCmd}) & 2\,d\\
4 & Geometry and BOM verification, report per item & 1\,d\\
5 & Vault: \code{import\_release} event and contract tests, provenance, attachments, fixed numbers in \code{run\_import} & 1.5\,d\\
6 & Migrate panel in FreeCAD (the import panel's pages: staging folder, CSV review, progress, report) & 1.5\,d\\
7 & Dry run on \code{LE-02001}, then all LE; FC-T01 checklist page; write-up & 1\,d\\
\midrule
 & Total & 10\,d\\
\bottomrule
\end{tabularx}
\end{center}
Steps 1--4 are useful before M6: the pilot (\eref{PDM-07}) can use them to bring the pedestal in if no FreeCAD model
exists.

\paragraph{Risks.} \emph{The Fusion API} changes between Fusion releases and runs only inside Fusion on your Mac; the
script is kept small and tested on one design first. \emph{Mesh and surface bodies} in a design have no solid to check;
they are listed and imported as they are. \emph{Large designs} make large STEP files that iCloud uploads slowly; the
staging folder sits outside iCloud. \emph{Name drift}: components renamed in Fusion between export runs map to new rows;
the CSV shows them side by side. \emph{The headless link-group crash} is avoided, not fixed.

%----------------------------------------------------------------------
\paragraph{Your answers to the questions below (24 Oct).} All eleven as recommended. 1: the nine FC-T01 criteria are
M6's checklist (criterion~8's two weeks can be shorter). 2: \textbf{LE first}, the other programs one by one with the
same tools. 3: \textbf{the Fusion API script}, which you run; the Fusion bridge only with your go-ahead. 4: the current
released revision and the current work; earlier letters as superseded records. 5: each design's \file{.f3d}
(\file{.f3z}) attached write-once. 6: the register's state is kept; working \code{A.}$n$ items arrive as WIP. 7:
\textbf{remodel on change}. 8: the Fusion PDFs attached; a drawing is redrawn when its part next changes. 9: scripted
parts skip Fusion. 10: a material table in the vault, checked through mass. 11: Fusion goes read-only when FC-T01 is
met, the subscription kept one more period. \textbf{The pilot starts fresh} (\S\ref{sec:pdm07-plan}, Q1): \code{LE-02001}
is modelled in FreeCAD and nothing is imported, so steps 1--4 are not used by the pilot; migration stays M6 work, for
the other LE items, the other programs and their history.

\subsubsection*{Questions before PDM-08 starts}

\begin{enumerate}
  \item \textbf{FC-T01.} Are the nine criteria of section~5 right? What is missing?\\
        \emph{Recommendation:} adopt them as M6's checklist; criterion~8's two weeks can be shorter if nothing
        sent you back to Fusion.
  \item \textbf{Scope.} Leonardo only, or every program in Fusion (OW, SI, FO, AS)?\\
        \emph{Recommendation:} LE first; the others program by program with the same tools once LE passes.
  \item \textbf{How to export.} A Fusion API script (exact part numbers, mass properties, transforms), or STEP and PDF
        exported by hand?\\
        \emph{Recommendation:} the script. You run it in Fusion; the Fusion bridge configured on this Mac could also run
        it, but only with your go-ahead, since Fusion is your live tool.
  \item \textbf{How much history.} Only the current revision of each item, or every released letter?\\
        \emph{Recommendation:} the current released revision and the current work; earlier letters as superseded records,
        with an \file{.f3d} attached only where Fusion still has that version.
  \item \textbf{Fusion archives.} Attach each design's \file{.f3d} (\file{.f3z} for assemblies) to its item?\\
        \emph{Recommendation:} yes, as write-once attachments; they are the only copy of the history.
  \item \textbf{Released on arrival.} Import released LE items as released (recorded as \emph{released by migration}),
        or as WIP needing a new review?\\
        \emph{Recommendation:} keep the register's state; working \code{A.}$n$ items arrive as WIP.
  \item \textbf{Remodelling.} Keep migrated parts as STEP solids, or remodel them natively?\\
        \emph{Recommendation:} remodel on change: a migrated part stays a solid until its first ECO, then is remodelled
        or direct-edited (\eref{MOD-04}).
  \item \textbf{Drawings.} Attach the Fusion PDF as the released drawing of a migrated revision, or redraw every
        drawing in TechDraw?\\
        \emph{Recommendation:} attach the PDFs; redraw a part's drawing when the part next changes.
  \item \textbf{Scripted parts} (build123d, CadQuery; ADR-0002). They never were Fusion parts. Same pipeline?\\
        \emph{Recommendation:} no Fusion step: the item holds the generated STEP converted to \file{.FCStd}, with the
        script and its commit attached; a change regenerates it.
  \item \textbf{Materials.} STEP carries colours, not materials. Map Fusion's material names to FreeCAD's library with
        a table in the vault?\\
        \emph{Recommendation:} yes; the geometry check then also compares mass, which catches a wrong mapping.
  \item \textbf{When Fusion goes read-only.} The day FC-T01 is met, or later?\\
        \emph{Recommendation:} when FC-T01 is met; keep the subscription one more renewal period as a fallback for
        reading the archives.
\end{enumerate}
